\begin{abstract}
We present a comprehensive geometric and topological analysis of magic squares through novel computational frameworks that reveal their intrinsic mathematical structure. By embedding magic squares as point clouds in $\mathbb{R}^3$, we discover a rich geometric landscape characterised by stratified manifold structure, distinctive topological features, and elegant covariance relationships. Our analysis of all 880 unique $4 \times 4$ magic squares demonstrates that they exhibit consistent geometric properties: convex hull volumes ranging from 58.0 to 97.3 cubic units (mean $84.1 \pm 7.1$), vertex counts from 8 to 16 (mean $12.2 \pm 1.4$), and angular uniformity measures of $0.569 \pm 0.000$. We establish a fundamental zero-covariance theorem showing that magic squares are precisely those arrangements where four specific covariances between spatial coordinates and entry values vanish simultaneously. Through Earth Mover's Distance analysis, we identify natural clustering with nearest-neighbour distances of $0.021 \pm 0.011$, revealing a structured metric space. Persistent homology analysis uncovers multi-scale topological features that distinguish magic squares from random arrangements. Our manifold analysis reveals 880 connected components with high clustering coefficient (1.000) and short average path length (1.00), indicating a highly structured geometric space. These findings establish magic squares as a paradigmatic example of discrete mathematical objects with rich continuous geometric structure, opening new avenues for analysing combinatorial problems through topological and geometric lenses. The methodology extends naturally to higher-order magic squares, magic cubes, and related combinatorial structures, with applications in discrete optimisation, coding theory, and quantum information processing.
\end{abstract} 